\subsection{薄壳平移与统一二阶项}\label{subsec:power-second}
本节仍固定 $0<\tau<n$，写 $\mu=\mu_{\tau,n}$、$K=K_{\tau,n}$，并定义
\[A_C=\inf_{\|u\|_\infty\le C^{-1}}\frac{\mu(Q(u,1))}{K},\qquad C>1.\]
\begin{lemma}[平移损失]\label{lem:power-shift}
当 $C\to\infty$ 时，
\[A_C=1-\frac{\tau(n-\tau)}2C^{-2}+O_{n,\tau}(C^{-3}).\tag{S.1}\]
\end{lemma}
\begin{proof}
层析恒等式给出
\[\mu(A)=(n-\tau)\int_0^\infty |A\cap Q(0,s)|s^{\tau-n-1}\,ds.\tag{S.2}\]
令 $\delta=C^{-1}$。每个坐标的区间交叠长度随平移的绝对值单调不增，
故 $A_C$ 在 $u=(\delta,\ldots,\delta)$ 取得。相应交叠方体的边长为
\[L_\delta(s)=\begin{cases}
2s,&0\le s\le1-\delta,\\
s+1-\delta,&1-\delta\le s\le1+\delta,\\
2,&s\ge1+\delta.
\end{cases}\]
中心与平移方体的交叠体积差只在上述薄壳内非零。令 $s=1+\delta v$，
在 $-1\le v\le1$ 上一致地有
\[(2\min\{1,s\})^n-L_\delta(s)^n
=n2^{n-1}\delta(1-|v|)+O_n(\delta^2).\]
因壳远离零，$s^{\tau-n-1}=1+O_{n,\tau}(\delta)$，而
$\int_{-1}^1(1-|v|)\,dv=1$。代入 (S.2) 得
\[\mu(Q(0,1))-\mu(Q(u,1))
=(n-\tau)n2^{n-1}\delta^2+O_{n,\tau}(\delta^3).
\]
除以 $K=n2^n/\tau$ 即得 (S.1)。这里未在权的内部奇点处求导。
\end{proof}

\begin{theorem}[统一二阶端点]\label{thm:power-second}
令
\[T_n=\min\left\{1,\frac{n+4-\sqrt{n^2+16}}2\right\}.
\]
若 $0<\tau<T_n$，则
\[C^{\rm out}_{c,w_{\tau,n}}(1-t)-1
=t^{\tau/n}+t^{2\tau/n}+o_{n,\tau}(t^{2\tau/n}).\tag{S.3}\]
其中 $T_1=(5-\sqrt{17})/2$、$T_2=3-\sqrt5$，且 $n\ge3$ 时 $T_n=1$。
\end{theorem}
\begin{proof}
写 $b=\tau/n$、$z=t^b$。定理 \ref{thm:op-power-n-leading} 的长类证明中，
将 $\mu(Q(x,r))\ge Kr^\tau(1-C^{-1})^\tau$ 改用
$\mu(Q(x,r))\ge Kr^\tau A_C$，其余实际半径上确界和共同核心的步骤不变，得到
\[C^{\rm out}_{c,w_{\tau,n}}(1-t)-1
\le\frac z{A_C-z}+O_{n,\tau}(tC^{n-\tau}).\tag{S.4}\]
此式在 $C>2$、$A_C>z$、$0<t<1/2$ 时成立，隐常数独立于输入及 $C$。
取 $C=t^{-a}$，由引理 \ref{lem:power-shift}，
\[\frac z{A_C-z}=z+z^2+O_{n,\tau}(z^3+zC^{-2}).\]
可同时取
\[\frac b2<a<\frac{1-2b}{n-\tau}\tag{S.5}\]
当且仅当 $\tau<(n+4-\sqrt{n^2+16})/2$。于是右端余项和短类均为 $o(z^2)$。
另一方面，前一定理的环形输入给出
\[C^{\rm out}_{c,w_{\tau,n}}(1-t)-1
\ge\frac{t^b}{(1+t^{1/n})^\tau-t^b}
=z+z^2+O_{n,\tau}(z^3+zt^{1/n}).\tag{S.6}\]
因 $\tau<1$，有 $zt^{1/n}=o(z^2)$，从而上下界匹配。
最后，$n\ge3$ 时 $\sqrt{n^2+16}\le n+2$，给出所述 $T_n$ 的化简。
\end{proof}
\begin{corollary}
在定理 \ref{thm:power-second} 的范围内，令 $a_*=\min_i a_i$，则共同中心多集合模型满足
\[C^{\boldsymbol a,\rm out}_{c,w_{\tau,n}}(1-\varepsilon)-1
=a_*^{-\tau/n}\varepsilon^{\tau/n}
+a_*^{-2\tau/n}\varepsilon^{2\tau/n}+o(\varepsilon^{2\tau/n}).\]
\end{corollary}
\begin{proof}
代入 $s_{\boldsymbol a}(1-\varepsilon)=\varepsilon/a_*+O(\varepsilon^2)$。
因 $\tau/n<1$，变量替换产生的 $O(\varepsilon^{1+\tau/n})$ 为第二尺度的小量。
\end{proof}
